\subsection{Взвешенный ICP для устойчивой регистрации}
\label{sec:2-3-icp-weighting}

Для регистрации соседних облаков точек в работе используется ICP с оценкой преобразования между парами точек с учётом весов. Архитектурно рассматриваемая реализация опирается на идеи KISS-ICP~\cite{kiss-icp}, однако адаптирует их под более простой сценарий обработки последовательности лидарных кадров. В статье KISS-ICP регистрация выполняется по схеме frame-to-map, то есть очередной кадр совмещается с локальной картой. В данной работе локальная карта в самом ICP не используется: преобразование оценивается непосредственно между двумя соседними облаками точек. При этом ключевые идеи статьи сохраняются: применяется двойная вокселизация, начальное приближение строится по предыдущей оценке движения, порог для поиска соответствий выбирается адаптивно, а для подавления выбросов используются веса, задаваемые функцией Гемана---Макклюра.

\subsubsection{Вокселизация и двойное прореживание облака точек}

Первый шаг подготовки облака точек к регистрации состоит в вокселизации. Во-первых, она уменьшает число точек и тем самым снижает вычислительную стоимость поиска соответствий. Во-вторых, она уменьшает избыточность в плотных областях сцены, где соседние точки несут почти одинаковую геометрическую информацию. В результате регистрация опирается не на все исходные измерения, а на более компактное представление сцены.

Как и в KISS-ICP~\cite{kiss-icp}, в реализации используется не одна, а две последовательные операции вокселизации. Сначала исходное облако прореживается с более мелким шагом; полученный набор точек обозначим через
\begin{equation}
    P_{\mathrm{merge}} = \mathcal{V}_{\alpha v}(P),
\end{equation}
где \(P\) --- исходное облако точек, \(v\) --- базовый размер вокселя, а \(\alpha\) --- коэффициент масштаба для первого прореживания. Затем выполняется второе, более грубое прореживание:
\begin{equation}
    \hat{P} = \mathcal{V}_{\beta v}(P_{\mathrm{merge}}),
\end{equation}
где \(\beta\) задаёт масштаб второго шага вокселизации. Именно облако \(\hat{P}\) используется в ICP, тогда как более плотное представление \(P_{\mathrm{merge}}\) сохраняет больше геометрической информации.

В программной реализации эти два шага выполняются последовательно в методе обработки кадра, как показано на рисунке~\ref{lst:icp-double-voxelization}. Сначала строится облако \texttt{points\_merge}, затем --- ещё более компактное облако \texttt{points\_reduced}. Идея полностью соответствует KISS-ICP, однако в данной работе второе облако используется для регистрации соседних кадров, а не для регистрации кадра с локальной картой.

\begin{figure}[h]
  \fontsize{12pt}{14pt}\selectfont
  \lstinputlisting[firstline=91,lastline=105,language=Python]{../src/lidar_odometry/actor.py}
  \caption{Двойная вокселизация облака точек перед регистрацией}
  \label{lst:icp-double-voxelization}
\end{figure}

\subsubsection{Начальное приближение по предыдущей оценке движения}

Следующая важная идея, заимствованная из KISS-ICP~\cite{kiss-icp}, состоит в том, что ICP не запускается из тождественного преобразования. Вместо этого используется кинематически осмысленное начальное приближение, построенное по предыдущей оценке движения. На малом интервале времени между соседними лидарными кадрами движение автомобиля обычно меняется плавно, поэтому модель постоянной скорости даёт разумный прогноз относительного смещения.

Если через \(\xi_{t-1}\) обозначить вектор линейных и угловых скоростей, оценённый на предыдущем шаге, а через \(\Delta t\) --- интервал времени между кадрами, то предсказанное относительное преобразование записывается как
\begin{equation}
    T_{\mathrm{pred}, t} = \exp(\xi_{t-1} \Delta t).
\end{equation}
В коде это соответствует построению объекта преобразования через экспоненциальные координаты с помощью функции SciPy~\cite{scipy-rigidtransform-from-exp-coords}; соответствующий фрагмент приведён на рисунке~\ref{lst:icp-initial-guess}. После этого текущее облако точек предварительно переносится в глобальную систему координат с учётом предыдущей позы и прогноза движения.

\begin{figure}[h]
  \fontsize{12pt}{14pt}\selectfont
  \lstinputlisting[firstline=107,lastline=110,language=Python]{../src/lidar_odometry/actor.py}
  \caption{Построение начального приближения по предыдущей оценке движения}
  \label{lst:icp-initial-guess}
\end{figure}

Такой шаг особенно важен в условиях перекрытия поля зрения. Если часть сцены оказывается скрыта динамическим объектом, то множество корректных соответствий уменьшается, и запуск ICP из нулевого смещения становится менее надёжным. Начальное приближение, построенное по предыдущему движению, сужает область поиска и повышает устойчивость последующей оценки преобразования.

\subsubsection{Адаптивный порог для поиска соответствий}

В классическом ICP после поиска ближайших соседей обычно вводится фиксированный порог \(\tau\), и пары точек с расстоянием больше этого порога отбрасываются. Недостаток такого подхода состоит в том, что один и тот же порог плохо переносится между различными режимами движения. При плавном движении слишком большой порог пропускает лишние выбросы, а при более резком манёвре слишком маленький порог может отбросить полезные соответствия.

В KISS-ICP~\cite{kiss-icp} предлагается выбирать этот порог адаптивно, исходя из того, насколько фактическое движение отклоняется от прогноза по модели постоянной скорости. Та же идея используется и в данной работе. Пусть \(\Delta T\) обозначает поправку, которую ICP вносит в предсказанное преобразование. Тогда верхняя оценка смещения точки записывается как сумма вкладов от поступательного и вращательного отклонений:
\begin{equation}
    \delta(\Delta T) = \delta_{\mathrm{rot}}(\Delta R) + \delta_{\mathrm{trans}}(\Delta t),
\end{equation}
где
\begin{equation}
    \delta_{\mathrm{trans}}(\Delta t) = \|\Delta t\|_2,
\end{equation}
\begin{equation}
    \delta_{\mathrm{rot}}(\Delta R) = 2 r_{\max} \sin\left(\frac{\theta}{2}\right).
\end{equation}
Здесь \(r_{\max}\) --- максимальная дальность учитываемых точек, а \(\theta\) --- угол вращения, соответствующий матрице \(\Delta R\). Именно такая формула используется в статье KISS-ICP для оценки максимально возможного смещения точки при отклонении от предсказанного движения.

В коде эта величина вычисляется в методе, показанном на рисунке~\ref{lst:icp-deviation}. Поступательная часть определяется как норма вектора переноса, а вращательная --- через угол поворота и максимальную дальность точки. Далее эти значения накапливаются по уже обработанным кадрам и используются для оценки текущего масштаба отклонений.

\begin{figure}[h]
  \fontsize{12pt}{14pt}\selectfont
  \lstinputlisting[firstline=173,lastline=194,language=Python]{../src/lidar_odometry/actor.py}
  \caption{Оценка отклонения от модели движения и вычисление адаптивного порога}
  \label{lst:icp-deviation}
\end{figure}

Если через \(M_t\) обозначить множество предыдущих шагов, на которых отклонение было достаточно заметным, то стандартное отклонение оценивается как
\begin{equation}
    \sigma_t = \sqrt{\frac{1}{|M_t|} \sum_{i \in M_t} \delta(\Delta T_i)^2},
\end{equation}
а сам порог выбирается по правилу трёх сигм:
\begin{equation}
    \tau_t = 3 \sigma_t.
\end{equation}
Именно эта логика реализована в методе \texttt{get\_deviation}: если статистика уже накоплена, то вычисляется \(\sigma_t\) и затем берётся \(\tau_t = 3\sigma_t\); если статистики ещё нет, используется начальное значение порога.

Таким образом, адаптивный порог в данной работе выбирается так же, как в статье KISS-ICP. Это важно для устойчивости алгоритма: при спокойном движении порог автоматически уменьшается и лучше отсекает выбросы, а при более резких изменениях движения увеличивается и не отбрасывает полезные пары точек слишком рано.

\subsubsection{Поиск соответствий между соседними облаками точек}

После построения начального приближения и выбора текущего порога \(\tau_t\) выполняется поиск соответствий между двумя соседними облаками точек. В отличие от KISS-ICP, где соответствия ищутся между текущим кадром и локальной картой, здесь они ищутся между текущим преобразованным облаком и предыдущим облаком, сохранённым в качестве опорного представления сцены. Тем самым сохраняется общая логика ICP, но упрощается структура данных и сама схема регистрации.

Для поиска ближайших соседей используется дерево \texttt{cKDTree}~\cite{scipy-ckdtree}. Для каждой точки источника выбирается ближайшая точка в целевом облаке, после чего применяется отсечение по адаптивному порогу:
\begin{equation}
    \|s_i - q_i\|_2 \leq \tau_t.
\end{equation}
Только пары, удовлетворяющие этому условию, участвуют в дальнейшей оценке преобразования. Соответствующий фрагмент кода приведён на рисунке~\ref{lst:icp-correspondences}.

\begin{figure}[h]
  \fontsize{12pt}{14pt}\selectfont
  \lstinputlisting[firstline=67,lastline=77,language=Python]{../src/lidar_odometry/icp.py}
  \caption{Поиск ближайших соответствий и отсечение по адаптивному порогу}
  \label{lst:icp-correspondences}
\end{figure}

\subsubsection{Оценка преобразования между парами точек с учётом весов}

После построения множества соответствий требуется оценить евклидово преобразование~\cite{scipy-rigidtransform}, которое наилучшим образом совмещает текущее облако с опорным. В данной работе, как и в KISS-ICP~\cite{kiss-icp}, используется point-to-point ICP. Однако оценка преобразования между парами точек выполняется не с одинаковым вкладом всех пар, а с учётом весов. Применяются веса, дающие устойчивость к выбросам, задаваемые функцией Гемана---Макклюра. Такой подход обеспечивает снижение влияния выбросов и повышает устойчивость оценки движения.

Математически на каждой итерации ICP решается задача
\begin{equation}
    \Delta T_{\mathrm{est}, j} = \arg\min_T \sum_{(s, q) \in C(\tau_t)} \rho\left(\|Ts - q\|_2\right),
\end{equation}
где \(C(\tau_t)\) --- множество соответствий, прошедших отсечение по порогу \(\tau_t\), а \(\rho\) --- функция Гемана---Макклюра. В статье KISS-ICP она записывается в виде
\begin{equation}
    \rho(e) = \frac{e^2 / 2}{\sigma_t / 3 + e^2}.
\end{equation}
В программной реализации напрямую вычисляются не значения функции стоимости, а соответствующие веса:
\begin{equation}
    w(e) = \frac{0.5}{\sigma_t / 3 + e^2}.
\end{equation}
Именно эта формула реализована в функции \texttt{geman\_mcclure\_weights}, показанной на рисунке~\ref{lst:icp-geman-mcclure}. Чем больше расстояние между соответствующими точками, тем меньше вес этой пары и тем слабее её влияние на итоговую оценку преобразования.

\begin{figure}[h]
  \fontsize{12pt}{14pt}\selectfont
  \lstinputlisting[firstline=17,lastline=19,language=Python]{../src/lidar_odometry/icp.py}
  \caption{Вычисление весов по функции Гемана---Макклюра}
  \label{lst:icp-geman-mcclure}
\end{figure}

После вычисления весов оценка преобразования между парами точек с учётом весов выполняется взвешенным вариантом алгоритма Кабша~\cite{kabsch-wikipedia,kabsch-nghiaho}. В коде для этого используется функция \texttt{Rotation.align\_vectors} из SciPy~\cite{scipy-align-vectors}, а затем по найденному вращению восстанавливается вектор переноса. Соответствующий фрагмент приведён на рисунке~\ref{lst:icp-kabsch}. Таким образом, задача минимизации и способ коррекции облака точек остаются теми же, что и в KISS-ICP: на каждой итерации вычисляется очередная поправка преобразования, после чего текущее облако дополнительно совмещается с опорным.

\begin{figure}[h]
  \fontsize{12pt}{14pt}\selectfont
  \lstinputlisting[firstline=35,lastline=47,language=Python]{../src/lidar_odometry/icp.py}
  \caption{Оценка преобразования между парами точек с учётом весов}
  \label{lst:icp-kabsch}
\end{figure}

\subsubsection{Итерационная схема ICP и критерий остановки}

Полный цикл ICP в реализации включает четыре повторяющихся шага: поиск соответствий, вычисление весов, оценку очередной поправки преобразования и обновление текущего положения облака точек. Если обозначить поправку на \(j\)-й итерации через \(\Delta T_j\), то накопленное преобразование обновляется по правилу
\begin{equation}
    T_{j+1} = \Delta T_j T_j.
\end{equation}
Именно такая композиция преобразований реализована в строке \texttt{transform = delta\_transform * transform}.

Критерий остановки также соответствует общей логике KISS-ICP: итерации прекращаются, когда очередная поправка становится достаточно малой. Если \(\xi(\Delta T)\) обозначает экспоненциальные координаты поправки, то условие остановки имеет вид
\begin{equation}
    \|\xi(\Delta T)\| < \gamma,
\end{equation}
где \(\gamma\) --- заданный порог сходимости. В коде это реализовано в функции \texttt{icp\_convergence\_criterion\_met}. При этом, в отличие от статьи, в реализации дополнительно сохраняется и верхняя граница на число итераций как защитный механизм от бесконечного цикла.

Полный фрагмент итерационной схемы ICP приведён на рисунке~\ref{lst:icp-main-loop}. В нём видно, что коррекция облака точек и задача минимизации организованы так же, как в KISS-ICP: сначала строятся соответствия, затем вычисляются веса, после чего находится очередная поправка преобразования и обновляется текущее положение источника. По своей математической структуре алгоритм сохраняет основные идеи KISS-ICP --- двойную вокселизацию, адаптивный порог, начальное приближение по предыдущему движению и point-to-point регистрацию, --- но в данной работе адаптирован к схеме регистрации соседних облаков точек без использования локальной карты внутри ICP.

\begin{figure}[h]
  \fontsize{12pt}{14pt}\selectfont
  \lstinputlisting[firstline=110,lastline=142,language=Python]{../src/lidar_odometry/icp.py}
  \caption{Основной цикл взвешенного ICP}
  \label{lst:icp-main-loop}
\end{figure}

